% GENERATED FILE. Edit canonical lessons, then run npm run generate:books.
% canonical-source-hash: fnv1a64:8d3c1f26c3493d30
% canonical-lessons: FA-C402-tahqiq, FA-C402-pazhuhesh, FA-C402-maqale, FA-C402-payan-name, FA-C402-sokhanrani

\chapter{Research and papers}
\label{ch:fa-research-and-papers}
\hlchaptermodality{402}

\begin{chapteropening}
\textbf{By the end of this chapter:} \emph{I can talk about research and studies, articles, a thesis and a lecture.}

Complete the last lesson of chapter 402: I can talk about research and studies, articles, a thesis and a lecture.
\end{chapteropening}

\section[\texorpdfstring{\fa{تحقیق} (tahqiq)}{تحقیق (tahqiq)}]{\fa{تحقیق} (tahqiq) — research, an investigation}

\label{lesson:FA-C402-tahqiq}



\begin{quote}
Before the new one: say the Persian for lecture notes, then the Persian for homework.
\end{quote}

\subsection*{You'll want to know: \fa{تحقیق}}
\textbf{\fa{تحقیق}} (\emph{tahqiq}) — \textquotedblleft{}research, an investigation\textquotedblright{}.

\begin{quote}
\fa{او} \fa{درباره} \fa{زبانهای} \fa{قدیمی} \fa{تحقیق} \fa{میکند}.
\end{quote}

\begin{quote}
\emph{u darbâre-ye zabânhâ-ye qadimi tahqiq mikonad.}
\end{quote}

\begin{quote}
She does research on ancient languages.
\end{quote}

\emph{Tahqiq kardan} — to research. Plural \emph{tahqiqât}.

\begin{cousinweb}
\textbf{\fa{تحقیق}} comes from Arabic \emph{taḥqīq} (EtymDB), a verbal noun from the same Arabic family as \emph{ḥaqīqa} \textquotedblleft{}truth\textquotedblright{}: research finds out what is true.
\end{cousinweb}

\begin{grammarlens}[title={What you've built}]
One more word for telling, explaining and describing at B1.
\end{grammarlens}

\subsection*{Your turn}
\begin{itemize}
\raggedright
  \item \emph{Say it:} \emph{tahqiq}
  \item \emph{Say it:} \emph{tahqiq}, once more
  \item \emph{Say it:} a sentence of your own with \emph{tahqiq}
  \item \emph{Recall:} say the Persian for lecture notes, then the Persian for homework, then say \emph{tahqiq} again
\end{itemize}

\begin{tcolorbox}[breakable,colback=teal!4,colframe=teal!35!black,title={Before you move on}]
What does \fa{تحقیق} mean? (\textquotedblleft{}Research, an investigation\textquotedblright{}.) Say it once more.
\end{tcolorbox}
\section[\texorpdfstring{\fa{پژوهش} (pazhuhesh)}{پژوهش (pazhuhesh)}]{\fa{پژوهش} (pazhuhesh) — scholarly research, a study}

\label{lesson:FA-C402-pazhuhesh}



\begin{quote}
Before the new one: say the Persian for homework, then the Persian for research.
\end{quote}

\subsection*{You'll want to know: \fa{پژوهش}}
\textbf{\fa{پژوهش}} (\emph{pazhuhesh}) — \textquotedblleft{}scholarly research, a study\textquotedblright{}.

\begin{quote}
\fa{این} \fa{پژوهش} \fa{پنج} \fa{سال} \fa{طول} \fa{کشید}.
\end{quote}

\begin{quote}
\emph{in pazhuhesh panj sâl tul keshid.}
\end{quote}

\begin{quote}
This study took five years.
\end{quote}

A native, formal near-synonym of \emph{tahqiq}. \emph{Pazhuheshgar} — a researcher.

\begin{cousinweb}
\textbf{\fa{پژوهش}} is native Persian, built on \emph{pazhuh-}, the stem of \emph{pazhuhidan} \textquotedblleft{}to search, to inquire\textquotedblright{}, with the noun suffix \emph{-esh}.
\end{cousinweb}

\begin{grammarlens}[title={What you've built}]
One more word for telling, explaining and describing at B1.
\end{grammarlens}

\subsection*{Your turn}
\begin{itemize}
\raggedright
  \item \emph{Say it:} \emph{pazhuhesh}
  \item \emph{Say it:} \emph{pazhuhesh}, once more
  \item \emph{Say it:} a sentence of your own with \emph{pazhuhesh}
  \item \emph{Recall:} say the Persian for homework, then the Persian for research, then say \emph{pazhuhesh} again
\end{itemize}

\begin{tcolorbox}[breakable,colback=teal!4,colframe=teal!35!black,title={Before you move on}]
What does \fa{پژوهش} mean? (\textquotedblleft{}Scholarly research, a study\textquotedblright{}.) Say it once more.
\end{tcolorbox}
\section[\texorpdfstring{\fa{مقاله} (maqâle)}{مقاله (maqâle)}]{\fa{مقاله} (maqâle) — an article, a paper}

\label{lesson:FA-C402-maqale}



\begin{quote}
Before the new one: say the Persian for research, then the Persian for scholarly research.
\end{quote}

\subsection*{You'll want to know: \fa{مقاله}}
\textbf{\fa{مقاله}} (\emph{maqâle}) — \textquotedblleft{}an article, a paper\textquotedblright{}.

\begin{quote}
\fa{او} \fa{یک} \fa{مقاله} \fa{درباره} \fa{آب} \fa{نوشت}.
\end{quote}

\begin{quote}
\emph{u yek maqâle darbâre-ye âb nevesht.}
\end{quote}

\begin{quote}
He wrote an article about water.
\end{quote}

Plural \emph{maqâlât} or \emph{maqâlehâ}.

\begin{cousinweb}
\textbf{\fa{مقاله}} was borrowed from Arabic \emph{maqāla} \textquotedblleft{}article, essay\textquotedblright{} (EtymDB).
\end{cousinweb}

\begin{grammarlens}[title={What you've built}]
One more word for telling, explaining and describing at B1.
\end{grammarlens}

\subsection*{Your turn}
\begin{itemize}
\raggedright
  \item \emph{Say it:} \emph{maqâle}
  \item \emph{Say it:} \emph{maqâle}, once more
  \item \emph{Say it:} a sentence of your own with \emph{maqâle}
  \item \emph{Recall:} say the Persian for research, then the Persian for scholarly research, then say \emph{maqâle} again
\end{itemize}

\begin{tcolorbox}[breakable,colback=teal!4,colframe=teal!35!black,title={Before you move on}]
What does \fa{مقاله} mean? (\textquotedblleft{}An article, a paper\textquotedblright{}.) Say it once more.
\end{tcolorbox}
\section[\texorpdfstring{\fa{پایان} \fa{نامه} (pâyân nâme)}{پایان نامه (pâyân nâme)}]{\fa{پایان} \fa{نامه} (pâyân nâme) — a thesis, a dissertation}

\label{lesson:FA-C402-payan-name}



\begin{quote}
Before the new one: say the Persian for scholarly research, then the Persian for an article.
\end{quote}

\subsection*{You'll want to know: \fa{پایان} \fa{نامه}}
\textbf{\fa{پایان} \fa{نامه}} (\emph{pâyân nâme}) — \textquotedblleft{}a thesis, a dissertation\textquotedblright{}.

\begin{quote}
\fa{خواهرم} \fa{پایان} \fa{نامه} \fa{اش} \fa{را} \fa{تمام} \fa{کرد}.
\end{quote}

\begin{quote}
\emph{khâharam pâyân nâme ash râ tamâm kard.}
\end{quote}

\begin{quote}
My sister finished her thesis.
\end{quote}

\emph{Pâyân} \textquotedblleft{}end\textquotedblright{} + \emph{nâme} \textquotedblleft{}letter, book\textquotedblright{}: the book at the end of a degree.

\begin{cousinweb}
\textbf{\fa{پایان} \fa{نامه}} is a native Persian compound built from \emph{pâyân} \textquotedblleft{}end\textquotedblright{} + \emph{nâme} \textquotedblleft{}letter, book\textquotedblright{}, a word EtymDB lists without a parent and usually traced to Middle Persian \emph{nāmag}. It is the book that ends a degree.
\end{cousinweb}

\begin{grammarlens}[title={What you've built}]
One more word for telling, explaining and describing at B1.
\end{grammarlens}

\subsection*{Your turn}
\begin{itemize}
\raggedright
  \item \emph{Say it:} \emph{pâyân nâme}
  \item \emph{Say it:} \emph{pâyân nâme}, once more
  \item \emph{Say it:} a sentence of your own with \emph{pâyân nâme}
  \item \emph{Recall:} say the Persian for scholarly research, then the Persian for an article, then say \emph{pâyân nâme} again
\end{itemize}

\begin{tcolorbox}[breakable,colback=teal!4,colframe=teal!35!black,title={Before you move on}]
What does \fa{پایان} \fa{نامه} mean? (\textquotedblleft{}A thesis, a dissertation\textquotedblright{}.) Say it once more.
\end{tcolorbox}
\section[\texorpdfstring{\fa{سخنرانی} (sokhanrâni)}{سخنرانی (sokhanrâni)}]{\fa{سخنرانی} (sokhanrâni) — a lecture, a talk}

\label{lesson:FA-C402-sokhanrani}



\begin{quote}
Before the new one: say the Persian for an article, then the Persian for a thesis.
\end{quote}

\subsection*{You'll want to know: \fa{سخنرانی}}
\textbf{\fa{سخنرانی}} (\emph{sokhanrâni}) — \textquotedblleft{}a lecture, a talk\textquotedblright{}.

\begin{quote}
\fa{سخنرانی} \fa{استاد} \fa{یک} \fa{ساعت} \fa{طول} \fa{کشید}.
\end{quote}

\begin{quote}
\emph{sokhanrâni-ye ostâd yek sâ'at tul keshid.}
\end{quote}

\begin{quote}
The professor's lecture lasted an hour.
\end{quote}

\emph{Sokhan} \textquotedblleft{}speech\textquotedblright{} + \emph{rân} \textquotedblleft{}driving\textquotedblright{}, from \emph{rândan}, + \emph{-i}. \emph{Sokhanrâni kardan} — to give a talk.

\begin{cousinweb}
\textbf{\fa{سخنرانی}} is native Persian, built from \emph{sokhan} \textquotedblleft{}speech, word\textquotedblright{} + \emph{rân-}, the present stem of \emph{rândan} \textquotedblleft{}to drive\textquotedblright{}, + \emph{-i}. A lecturer drives the speech forward.
\end{cousinweb}

\begin{grammarlens}[title={What you've built}]
One more word for telling, explaining and describing at B1.
\end{grammarlens}

\subsection*{Your turn}
\begin{itemize}
\raggedright
  \item \emph{Say it:} \emph{sokhanrâni}
  \item \emph{Say it:} \emph{sokhanrâni}, once more
  \item \emph{Say it:} a sentence of your own with \emph{sokhanrâni}
  \item \emph{Recall:} say the Persian for an article, then the Persian for a thesis, then say \emph{sokhanrâni} again
\end{itemize}

\begin{tcolorbox}[breakable,colback=teal!4,colframe=teal!35!black,title={Before you move on}]
What does \fa{سخنرانی} mean? (\textquotedblleft{}A lecture, a talk\textquotedblright{}.) Say it once more.
\end{tcolorbox}
